% FILE: 07_open_questions_and_next_work.tex
% Project: ggen — Governance Generation Engine

\section{Open Questions and Next Work}
\label{sec:ggen-open-questions}

\subsection{Known Gaps}
\label{sec:ggen-open-questions-gaps}

Four structural gaps are identified in the available artifacts.

\textbf{Gap 1: ggen binary non-existence.} The \texttt{src/} directory is empty. ggen is
declared as a manufacturing engine with a command-line interface (\texttt{ggen --config
ggen.toml --execute-all}), but no compiled binary or Rust source file implementing that interface
exists in the repository. Whether ggen is an external tool installed in the environment, a cargo
plugin, or whether the manufacturing runs were performed by manually authoring outputs against
the template contract is not documented in any available artifact. The BLAKE3 receipt chain
proves that manufacturing operations occurred (timestamps, input hashes, output hashes, and
signatures are present), but the agent that performed those operations cannot be identified from
the repository surface alone.

\textbf{Gap 2: Visualizer generation rule disabled.} The
\texttt{visualizer-dashboard-nextjs} generation rule is commented out in \texttt{ggen.toml} with
the note ``DISABLED: visualizer-dashboard-nextjs template parse error - deferred to Phase 6
recovery.'' The template file \texttt{templates/visualizer-dashboard.tsx.tera} would have
manufactured a NextJS page at
\texttt{../experiments/visualizer-nextjs/src/app/page.tsx}. The nature of the Tera template
parse error and the Phase 6 recovery plan are not documented in any available artifact.

\textbf{Gap 3: Empty output hashes in receipts 2 and 3.} The chain-root receipt
(\texttt{sync-20260601-175316.json}) records the BLAKE3 hash of \texttt{blue\_river\_dam/src/lib.rs}
as a concrete output. The second receipt (\texttt{sync-20260601-203853.json}) and third receipt
(\texttt{sync-20260601-230600.json}) both record \texttt{output\_hashes: []}. This means the
second and third manufacturing runs committed receipts without producing or hashing new output
artifacts. Whether this reflects a deliberate re-run of the configuration step only, or an
execution failure that the receipt chain does not surface, is unresolved.

\textbf{Gap 4: Ontology population uncertainty.} The SPARQL query endpoint is configured as
\texttt{in-memory} in \texttt{ggen.toml}. The board admissibility query
(\texttt{extract-board-claims.rq}) requires \texttt{wasm4pm:ConformanceVerdict} instances with
fitness $\geq 0.95$ and precision $\geq 0.90$ to be present in the in-memory triple store at
manufacturing time. The TTL files (\texttt{ontology-extensions.ttl},
\texttt{wasm4pm-compat.ttl}) define the schema but do not contain pre-populated verdict
instances. Whether the board deck and diligence workbook were ever manufactured with actual
conformance data, or whether those generation rules were not activated because the in-memory
store was empty at manufacturing time, is not documented.

\subsection{Deferred Gates}
\label{sec:ggen-open-questions-deferred}

The following proof gates are declared in the manufacturing summary but not confirmed as having
passed:

\begin{itemize}
  \item \textbf{Audit gate execution}: Seven audit scripts covering 42 or more named gates exist
    as manufactured artifacts. Whether these scripts have been executed against the actual
    wasm4pm-compat Rust source and produced passing results is not recorded in the BLAKE3
    receipt chain. The \texttt{audit.json} records \texttt{validation\_passed: true} but its
    \texttt{inputs}, \texttt{pipeline}, and \texttt{outputs} fields are all empty, suggesting
    the validation was structural (schema check) rather than executable (gate execution).
  \item \textbf{Test execution against compiled artifact}: The embedded test suites in
    \texttt{witness-markers-20260601.rs} and \texttt{blue\_river\_dam/src/lib.rs} declare lawful
    algebraic proof obligations. Whether \texttt{cargo test} has been run against the compiled
    manufactured artifacts and all tests have passed is not documented in the receipt chain.
  \item \textbf{PROCESS\_INTELLIGENCE\_ALIVE\_001}: The checkpoint is referenced as pending in
    \texttt{GGEN\_MANUFACTURING\_SUMMARY.md}. No committed checkpoint file for this verdict
    has been located.
\end{itemize}

\subsection{Research Risks}
\label{sec:ggen-open-questions-risks}

\textbf{Risk 1: Circular provenance.} The \texttt{[[inference.rules]]} normalize rule uses
\texttt{FILTER(false)}, making it a no-op. If meaningful inference was intended to populate
conformance verdict instances in the in-memory triple store before executing the board claims
query, the no-op means no inferred facts were ever produced. In this scenario, the manufactured
M\&A deck would be an empty artifact (no claims passing the FILTER) rather than a substantive
board deliverable. This risk cannot be resolved without evidence of what RDF facts were loaded
into the in-memory store at manufacturing time.

\textbf{Risk 2: Template-only manufacturing.} The GGEN\_MANUFACTURING\_SUMMARY.md's Phase 3
scope is the conversion of manifests into governance law files, templates, and audit scripts.
This is a form of manufacturing---specification artifacts manufactured from higher-order
specifications---but it is one level of abstraction removed from the pipeline's declared purpose
of manufacturing capital-market artifacts from conformance verdicts. The research risk is that
ggen is documented as a manufacturing engine but has, to date, primarily manufactured its own
manufacturing specification rather than the downstream capital-market artifacts the specification
describes.

\textbf{Risk 3: Binary-external dependency.} If ggen is an external binary, it introduces an
undocumented dependency on a tool not versioned in this repository. Any future re-execution of
the manufacturing pipeline requires that external tool to be available at the same version. The
receipt chain records \texttt{ggen\_version: 26.5.21} but does not record where or how to obtain
that binary, creating a reproducibility risk.

\subsection{Next Receipted Work}
\label{sec:ggen-open-questions-next}

The following work items, each requiring a new BLAKE3 receipt to be considered complete, are
indicated by the gap analysis:

\begin{enumerate}
  \item \textbf{Phase 6 recovery: visualizer-dashboard-nextjs.} Diagnose the Tera template
    parse error in \texttt{templates/visualizer-dashboard.tsx.tera}, repair the template, and
    re-enable the generation rule in \texttt{ggen.toml}. Execute the rule, confirm the NextJS
    artifact is manufactured at
    \texttt{../experiments/visualizer-nextjs/src/app/page.tsx}, and seal with a BLAKE3 receipt
    recording a non-empty \texttt{output\_hashes}.
  \item \textbf{Ontology population and board deck manufacturing.} Populate the in-memory triple
    store with actual \texttt{wasm4pm:ConformanceVerdict} instances (sourced from wasm4pm
    execution outputs) before executing \texttt{extract-board-claims.rq}. Execute the board deck
    and diligence workbook generation rules and seal with BLAKE3 receipts recording non-empty
    \texttt{output\_hashes}.
  \item \textbf{Audit gate execution and receipt.} Execute all seven audit scripts against the
    actual wasm4pm-compat source. Record pass/fail results and seal the audit execution log with
    a BLAKE3 receipt. Update \texttt{audit.json} to record non-empty \texttt{inputs},
    \texttt{pipeline}, and \texttt{outputs} fields.
  \item \textbf{Test execution and receipt.} Run \texttt{cargo test} against the compiled
    \texttt{blue\_river\_dam} and \texttt{wasm4pm-compat} crates. Record the test output
    (pass count, test names) and seal with a BLAKE3 receipt referencing the compiled artifact
    hashes.
  \item \textbf{PROCESS\_INTELLIGENCE\_ALIVE\_001 checkpoint.} Once all four preceding items
    are receipted and all gates pass, commit the \texttt{PROCESS\_INTELLIGENCE\_ALIVE\_001}
    checkpoint file to the receipts directory and update the ggen status from PARTIAL to ALIVE.
\end{enumerate}
